Large production clusters increasingly rely on elastic resource control to balance service quality against resource efficiency. In practice, many production stacks still prefer reactive threshold rules because they are operationally simple and easy to explain \citep{lorido2014review,Burns2016BorgKubernetes,Bernstein2014ContainersCloud}. However, reactive autoscaling often expands capacity only after demand has already surged and contracts capacity only after long idle periods. The result is a recurring trade-off between service-level objective violations and persistent over-provisioning.

Predictive autoscaling is a natural response to this limitation, but the literature still shows two recurring weaknesses. First, many papers evaluate forecasting quality and autoscaling quality separately, which makes it difficult to judge how improved prediction changes control trade-offs in practice. Second, many strong cluster-management studies depend on private traces or platform-specific deployments that are hard to reproduce. This reproducibility gap weakens both empirical comparability and operational credibility.

This problem setting is already visible in prior work on host-load prediction \citep{Dinda2000HostLoad}, container autoscaling \citep{Casalicchio2019ContainerAutoscaling}, and uncertainty-aware workload forecasting with Alibaba and Google traces \citep{Rossi2025UncertaintyAwareForecasting}. Taken together, these studies show that forecasting, autoscaling, and distributed resource management are closely linked, but they also show that forecasting quality, control behavior, and validity threats need to be analyzed jointly rather than in isolation.

The present paper therefore asks a narrower question: can public cluster traces support a reproducible predictive capacity-control evaluation that clarifies when decision trade-offs improve relative to reactive rules without requiring private infrastructure? The study instantiates this question on the Alibaba Cluster Trace Program \citep{alibabaclusterdata2026}, using the \texttt{cluster-trace-v2018} release whose official documentation reports about 4000 machines observed over 8 days \citep{alibabav2018trace}. Public analyses of Alibaba traces have already shown rich workload diversity and nontrivial resource-efficiency structure \citep{Jiang2022Characterizing,guo2019limits,luo2021characterizing}, making the dataset suitable for a trace-driven systems study.

The forecasting layer is coupled explicitly to the control layer. Instead of treating prediction as an isolated machine-learning benchmark, the experimental pipeline routes multi-horizon forecasts into an autoscaling simulator with headroom, cooldown, and safety-margin constraints. This separation keeps the control logic interpretable and allows the forecasting backbone to be replaced later without changing the rest of the evaluation route.

The evidence base is also intentionally broader than a single aggregate benchmark. In addition to the aggregate cluster series, the study includes a 12-machine high-coverage cohort, a cross-machine transfer benchmark, and a rule-based broad service-level cohort extracted from official \texttt{container\_meta} and a full mirrored pass over \texttt{container\_usage}. Aggregate and machine-level results show that simple linear baselines remain strong point predictors, whereas \texttt{UA-MSTCN-Lite} contributes calibrated quantiles that are useful for predictive control. The service-level evidence narrows the construct-validity gap further by testing the same forecasting-to-policy route at an application-level autoscaling unit.

The present implementation is intentionally conservative in what it claims. The forecasting backbone is a lightweight uncertainty-aware surrogate rather than the final deep UA-MSTCN architecture, and the service-level cohort is a rule-based broad benchmark rather than a full service census. These limits are kept explicit so that the results are interpreted as verified evidence for the current implementation, not as unsupported claims about a deeper architecture or a complete production deployment. The contribution is therefore evaluative: a reproducible public-trace benchmark and boundary-mapping study rather than a claim that the current controller or surrogate is universally superior.

This paper makes four evaluation-oriented contributions:
\begin{enumerate}
  \item It formulates predictive cluster capacity control as a reproducible public-trace evaluation problem rather than a private platform study.
  \item It establishes a public-trace benchmark pipeline from Alibaba trace acquisition through aggregate-, machine-, and service-level preprocessing, forecasting, and trace-driven policy simulation.
  \item It contributes a multi-granularity evidence surface spanning aggregate forecasting, machine-level benchmarking, cross-machine transfer, and service-level heterogeneity under one consistent protocol.
  \item It reports a boundary-mapping result for the current predictive route: simple linear baselines remain strong for point prediction, uncertainty-aware quantiles remain useful for control, and broader audited service coverage shows that the current shared global predictive policy is not validated as a service-level improvement route.
\end{enumerate}

The remainder of the paper is organized as follows. Section~2 reviews adjacent work on workload characterization, forecasting, autoscaling, and generalization. Section~3 formalizes the forecasting and control problem. Section~4 describes the dataset routes, preprocessing, models, and evaluation protocol. Section~5 reports the forecasting and policy results. Section~6 discusses validity threats, and Section~7 concludes the paper.
